% ============================================================
%  Formelsammlung INF4/2 (Low level & seminumerische Algorithmen)
%  Neue Kapitel als \input{parts/XX_name} einbinden.
%  Layout-Template übernommen aus SOS_B_FS.
% ============================================================
\documentclass[a4paper,10pt]{article}
\usepackage[utf8]{inputenc}
\usepackage[T1]{fontenc}
\usepackage[ngerman]{babel}
\usepackage{amsmath,amssymb,mathtools}
\usepackage{xcolor}
\usepackage{geometry}
\usepackage{array}
\usepackage{microtype}
\usepackage{needspace}
\usepackage{fancyhdr}
\geometry{top=8mm,bottom=4mm,left=10mm,right=10mm,includehead,includefoot}
\setlength{\headheight}{11pt}\setlength{\headsep}{1.2mm}\setlength{\footskip}{3.4mm}
\setlength{\parskip}{0pt}
\setlength{\parindent}{0pt}
\renewcommand{\baselinestretch}{0.9}
\setlength{\emergencystretch}{2em}
\hfuzz=3pt
\setlength{\fboxsep}{0.18em}
\setlength{\fboxrule}{0.3pt}

% ---- Farbschema: eine Farbe pro Teil ---------------------
\colorlet{cNewton}{blue}            % 1 Newton & Fixpunkte (Klausurschwerpunkt)
\colorlet{cBits}{teal}              % 2 Bit Wizardry
\colorlet{cOpt}{orange!90!red}      % 3 Code-Optimierung
\colorlet{cRel}{violet}             % 4 Asymptotik & Relationen

% ---- Teil-Kopf -------------------------------------------
\newcounter{bereichnr}
\newcommand{\teilnr}{}
\newcommand{\teilname}{}
\newcommand{\teilkopf}[2]{%
  \renewcommand{\teilnr}{#1}%
  \renewcommand{\teilname}{#2}%
  \setcounter{bereichnr}{0}%
  \markboth{#2}{#2}%
}

% ---- Bereichskopf -----------------------------------------
\newcommand{\bereich}[2][cNewton]{%
  \needspace{4\baselineskip}%
  \refstepcounter{bereichnr}%
  \markboth{\teilname}{\teilname}%
  \vspace{0.5em}%
  \noindent\colorbox{#1!10}{%
    \makebox[\dimexpr\linewidth-2\fboxsep\relax][l]{%
      \small\bfseries\color{#1!80!black}\strut\enspace
      \colorbox{#1!25}{\textcolor{#1!75!black}{\,\teilnr.\thebereichnr\,}}\enspace #2}}%
  \nopagebreak%
  \vspace{0.1em}\par\noindent%
}

\newcommand{\bereichende}{%
  \vspace{0.2em}%
  {\color{gray!50}\hrule height 0.3pt}%
  \vspace{0.2em}%
}

% ---- Drei-Spalten-Layout ----------------------------------
\newenvironment{formelreihe}{%
  \noindent\footnotesize%
  \setlength{\tabcolsep}{0pt}%
  \renewcommand{\arraystretch}{0.85}%
  \begin{tabular}{@{}p{0.32\linewidth}@{\hspace{0.02\linewidth}}%
                     p{0.32\linewidth}@{\hspace{0.02\linewidth}}%
                     p{0.32\linewidth}@{}}%
}{%
  \end{tabular}\par%
}

\newcommand{\formelblock}[2]{%
  \parbox[t]{\linewidth}{%
    {\scriptsize\scshape\color{gray!65!black} #1}\\[0.03em]%
    $\displaystyle #2$%
  }%
}

\newcommand{\beispielblock}[2]{%
  \fcolorbox{orange!50}{yellow!15}{%
    \parbox[t]{\dimexpr\linewidth-2\fboxsep-2\fboxrule\relax}{%
      {\scriptsize\scshape\color{orange!80!black} Bsp: #1}\\[0.03em]%
      $\displaystyle #2$%
    }}%
}

\newcommand{\infoblock}[1]{%
  \parbox[t]{\linewidth}{\scriptsize\itshape\color{black!60} #1}%
}

% ---- Rezept ------------------------------------------------
\newcounter{rezeptnr}
\newenvironment{rezept}{%
  \par\nobreak\noindent
  \setcounter{rezeptnr}{0}%
  \footnotesize\setlength{\tabcolsep}{2pt}%
  \renewcommand{\arraystretch}{1.05}%
  \begin{tabular}{@{}r@{\;\,}p{\dimexpr\linewidth-1.6em\relax}@{}}%
}{%
  \end{tabular}\par%
}
\newcommand{\rs}[1]{\refstepcounter{rezeptnr}\therezeptnr.&#1\\}

\newcommand{\fallstrick}[1]{%
  \par\vspace{0.1em}\noindent
  {\footnotesize\color{red!72!black}\textbf{Fallstricke:}}\;{\footnotesize #1}\par}

\newcommand{\miniexample}[2]{%
  \par\vspace{0.12em}%
  \fcolorbox{orange!50}{yellow!15}{%
    \parbox{\dimexpr\linewidth-2\fboxsep-2\fboxrule\relax}{%
      {\footnotesize\scshape\color{orange!80!black} Bsp: #1}\\[0.12em]%
      \footnotesize #2}}\par}

% ---- Code-Zeile kompakt, math- und textmode-tauglich -------
\newcommand{\code}[1]{\text{\ttfamily\footnotesize #1}}

% ---- Kopf-/Fußzeile ----------------------------------------
\pagestyle{fancy}
\fancyhf{}
\fancyhead[L]{\small\bfseries Formelsammlung INF4/2}
\fancyhead[R]{\small\color{gray!45!black}\rightmark}
\fancyfoot[C]{\small\color{gray!55!black}Seite \thepage}
\renewcommand{\headrulewidth}{0.4pt}
\renewcommand{\footrulewidth}{0pt}

% ==========================================================
\begin{document}

% ============================================================
%  Teil 1 — Newton & Fixpunkte  (Klausurschwerpunkt, SS16-bestätigt)
% ============================================================
\teilkopf{1}{Newton \& Fixpunkte}

\bereich[cNewton]{Fixpunkte \& Typen}
\begin{formelreihe}
\formelblock{Iteration / Fixpunkt}{x_{k+1}=\Phi(x_k),\quad \Phi(x_*)=x_*}
&
\formelblock{Fehler-Mechanik (Taylor)}{\Phi(x_*+\varepsilon)\approx x_*+\Phi'(x_*)\,\varepsilon}
&
\beispielblock{$\Phi=x^2$}{\begin{aligned}[t]&x_*{=}0{:}\ \Phi'{=}0\ \text{superattr.}\\ &x_*{=}1{:}\ \Phi'{=}2\ \text{repulsiv}\end{aligned}}
\\[0.4em]
\formelblock{attraktiv}{|\Phi'(x_*)|<1\ \Rightarrow\ \text{Konvergenz}}
&
\formelblock{repulsiv}{|\Phi'(x_*)|>1\ \Rightarrow\ \text{Divergenz}}
&
\formelblock{superattraktiv}{\Phi'(x_*)=0\ \Rightarrow\ \text{Ordnung}\ \geq 2}
\end{formelreihe}
\infoblock{Geometrie: FP = Schnitt $y{=}\Phi(x)$ mit Diagonale $y{=}x$. Basin der Attraktion = alle konvergenten $x_0$; Rand i.\,Allg.\ fraktal.}
\bereichende

\bereich[cNewton]{Konvergenzordnung — das Klausurschema}
\begin{formelreihe}
\formelblock{Def.\ (rel.\ Fehler, $x_*\neq 0$)}{\Phi\bigl(x_*(1{+}\varepsilon)\bigr)=x_*\bigl(1+\alpha\,\varepsilon^{n}+O(\varepsilon^{n+1})\bigr)}
&
\formelblock{äquivalent}{\Phi'(x_*){=}\dots{=}\Phi^{(n-1)}(x_*){=}0,\ \Phi^{(n)}(x_*){\neq}0}
&
\beispielblock{ablesen}{\begin{aligned}[t]&\Phi{=}x_*(1-\tfrac{3}{2}\varepsilon^2-\tfrac{1}{2}\varepsilon^3)\\ &\alpha_1{=}0,\ \alpha_2{\neq}0\ \Rightarrow\ \text{quadr.}\end{aligned}}
\end{formelreihe}
\begin{rezept}
\rs{$x=x_*(1+\varepsilon)$ in $\Phi$ einsetzen}
\rs{ausmultiplizieren, nach $\varepsilon$-Potenzen ordnen}
\rs{kleinster nicht verschwindender Koeff.\ $\Rightarrow$ Ordnung $n$}
\end{rezept}
\infoblock{Ordnung 1 ($|\alpha|{<}1$): konst.\ viele Stellen/Schritt $\cdot$ Ordnung 2: Stellen \textbf{verdoppeln} sich $\cdot$ Ordnung 3: verdreifachen.}
\fallstrick{Schema gilt nur für $x_*\neq 0$! Bei $x_*=0$ absoluten Fehler nehmen: $\varepsilon_{k+1}=\alpha\,\varepsilon_k^n$ (Bsp.\ $\Phi{=}x^2$: $\varepsilon_{k+1}{=}\varepsilon_k^2$).}
\bereichende

\bereich[cNewton]{Newton-Verfahren}
\begin{formelreihe}
\formelblock{Newton-Iteration für $f(r){=}0$}{N(x)=x-\tfrac{f(x)}{f'(x)}}
&
\formelblock{FP-Check}{N(r)=r-\tfrac{0}{f'(r)}=r\ \checkmark}
&
\beispielblock{NR-Helfer}{\bigl(\tfrac{f}{f'}\bigr)'=1-\tfrac{f\,f''}{(f')^2},\quad (x^{-2})'=-2x^{-3}}
\\[0.4em]
\formelblock{Superattraktivität allg.}{N'(x)=\tfrac{f(x)\,f''(x)}{f'(x)^2}\ \Rightarrow\ N'(r)=\tfrac{0\cdot f''(r)}{f'(r)^2}=0}
&
\formelblock{Voraussetzung}{f'(r)\neq 0\ \text{(einfache NS)}\ \Rightarrow\ \text{quadr.\ konv.}}
&
\beispielblock{Geometrie}{\begin{aligned}[t]&\text{Tangente: } y=f(x_k)+f'(x_k)(x-x_k)\\ &y{=}0 \Rightarrow x = N(x_k)\end{aligned}}
\end{formelreihe}
\bereichende

\bereich[cNewton]{Die drei Klausur-Klassiker}
\begin{formelreihe}
\formelblock{$1/d$ (Division verboten)}{\begin{aligned}[t]f&=\tfrac1x-d,\ f'=-\tfrac{1}{x^2}\\ N&=x(2-dx)=x+x(1-dx)\end{aligned}}
&
\formelblock{$1/\sqrt d$ (Div.\ \& Wurzel verboten)}{\begin{aligned}[t]f&=\tfrac{1}{x^2}-d,\ f'=-\tfrac{2}{x^3}\\ N&=\tfrac{x}{2}(3-dx^2)=x+x\,\tfrac{1-dx^2}{2}\end{aligned}}
&
\beispielblock{$\varepsilon$-Nachweis $1/d$}{\begin{aligned}[t]&N\bigl(\tfrac1d(1{+}\varepsilon)\bigr)=\tfrac1d(1{+}\varepsilon)(2-1-\varepsilon)\\ &=\tfrac1d(1{+}\varepsilon)(1{-}\varepsilon)=\tfrac1d(1-\varepsilon^2)\ \Rightarrow\ \text{quadr.}\end{aligned}}
\\[0.4em]
\formelblock{$\sqrt d$ direkt (enthält Division!)}{\begin{aligned}[t]f&=x^2-d,\ f'=2x\\ N&=\tfrac12\bigl(x+\tfrac dx\bigr)\end{aligned}}
&
\formelblock{$\sqrt d$ ohne Division}{\sqrt d = d\cdot\tfrac{1}{\sqrt d}\quad\text{(rechte Spalte oben)}}
&
\beispielblock{$1/\sqrt d$: FP \& SA}{\begin{aligned}[t]&N'(x)=\tfrac32(1-dx^2)\\ &N'(\pm\tfrac{1}{\sqrt d})=\tfrac32(1-1)=0\ \text{(SA)}\end{aligned}}
\\[0.4em]
\formelblock{$1/\sqrt d$: $\varepsilon$-Entwicklung}{N\bigl(\tfrac{1}{\sqrt d}(1{+}\varepsilon)\bigr)=\tfrac{1}{\sqrt d}\bigl(1-\tfrac32\varepsilon^2-\tfrac12\varepsilon^3\bigr)}
&
\formelblock{allg.\ Schema $d^{-1/a}$}{\begin{aligned}[t]&\Phi_2(x)=x\bigl(1+\tfrac{1-x^a d}{a}\bigr)\\ &\beta=-\tfrac{a+1}{2}\ \text{(quadr.)}\end{aligned}}
&
\beispielblock{unerwünschte FP von $N{=}x(2{-}dx)$}{\begin{aligned}[t]&N(0)=0\ \text{FP, }N'(0){=}2\ \text{repulsiv (gut!)}\\ &N(-\tfrac1d)=-\tfrac3d\neq-\tfrac1d\ \text{kein FP}\end{aligned}}
\end{formelreihe}
\begin{rezept}
\rs{NS formulieren: $f$ wählen mit $f(\text{gesucht})=0$; verbotene Op.\ darf in $N$ nicht auftauchen}
\rs{$f'$ bilden, $N=x-f/f'$ aufstellen, vereinfachen}
\rs{FP prüfen: $N(x_*)=x_*$ einsetzen}
\rs{SA zeigen: $N'(x_*)=0$ \textbf{oder} $\varepsilon$-Entwicklung ($\alpha_1=0$)}
\rs{ggf.\ weitere FP testen (0? $-x_*$?) und klassifizieren}
\end{rezept}
\fallstrick{$f=x-\tfrac1d$ ist sinnlos ($N\equiv\tfrac1d$ enthält Ergebnis) $\cdot$ $\sqrt d$ direkt enthält $\tfrac dx$ $\Rightarrow$ Division $\cdot$ Halbieren ($\tfrac x2$) ist erlaubt = Bit-Shift $\cdot$ SS16 fragt explizit FP bei $0$ und $-1/d$!}
\bereichende

% ============================================================
%  Teil 2 — Bit Wizardry
% ============================================================
\teilkopf{2}{Bit Wizardry}

\bereich[cBits]{Grundtricks (alles O(1))}
\begin{formelreihe}
\formelblock{Zweierkomplement}{-w=\mathord{\sim}w+1}
&
\formelblock{niedrigstes gesetztes Bit}{\begin{aligned}[t]&\code{w \& -w}\ \text{isolieren}\\ &\code{w \& (w-1)}\ \text{löschen}\end{aligned}}
&
\beispielblock{$w{=}01011000$}{\begin{aligned}[t]&\mathord{\sim}w{=}10100111,\ -w{=}10101000\\ &w\,\&\,{-w}=00001000\end{aligned}}
\\[0.4em]
\formelblock{Zweierpotenz-Test}{\code{!(x \& (x-1))}\ \text{(true auch bei 0)}}
&
\formelblock{Rotation rechts (BPL $=$ Wortbreite)}{\code{(w>{}>s) | (w<{}<(BPL-s))}}
&
\beispielblock{BPL$=$6, $w{=}$\code{abcdef}, $s{=}2$}{\begin{aligned}[t]&\code{w>{}>2}=\code{00abcd}\\ &\code{w<{}<4}=\code{ef0000}\ \Rightarrow\ \code{efabcd}\end{aligned}}
\end{formelreihe}
\infoblock{1-Cycle-Befehle: XOR, AND, NOT, ADD, SUB, NEG, INC, DEC, SHIFT, ROTATE $\cdot$ mit \code{unsigned long} arbeiten (signed shift zieht Vorzeichenbit nach).}
\bereichende

\bereich[cBits]{popcount (ones\_count) — Divide \& Conquer}
\begin{formelreihe}
\formelblock{Masken-Schema (Hex-Ziffern 5, 3, f)}{\begin{aligned}[t]&k{=}1{:}\ \code{0x5555\dots}\ (0101\dots)\\ &k{=}2{:}\ \code{0x3333\dots}\ (0011\dots)\\ &k{=}4{:}\ \code{0x0f0f\dots}\\ &k{=}8{:}\ \code{0x00ff\dots},\ k{=}16{:}\ \code{0x0000ffff}\end{aligned}}
&
\formelblock{Schritt (Blöcke $k\to 2k$)}{\code{x = (M \& x) + (M \& (x>{}>k))}}
&
\beispielblock{4 Bit, $x{=}1011$}{\begin{aligned}[t]&(0101\,\&\,x)+(0101\,\&\,x{>{}>}1)\\ &=0001+0101\to 0110\ \text{(je 2er-Block)}\\ &(x{>{}>}2)+(x\,\&\,0011)=01+10=3\ \checkmark\end{aligned}}
\end{formelreihe}
\miniexample{64-Bit-Referenz (fxtbook §1.8) — je Zeile Maske vor \& nach Shift}{%
\code{x = (0x5555555555555555 \& x) + (0x5555555555555555 \& (x>{}>~1));}\\
\code{x = (0x3333333333333333 \& x) + (0x3333333333333333 \& (x>{}>~2));}\\
\code{x = (0x0f0f0f0f0f0f0f0f \& x) + (0x0f0f0f0f0f0f0f0f \& (x>{}>~4));}\\
\code{x = (0x00ff00ff00ff00ff \& x) + (0x00ff00ff00ff00ff \& (x>{}>~8));}\\
\code{x = (0x0000ffff0000ffff \& x) + (0x0000ffff0000ffff \& (x>{}>16));}\\
\code{x = (0x00000000ffffffff \& x) + (0x00000000ffffffff \& (x>{}>32));}}
\infoblock{32 $\to$ 64 Bit: (1) Masken verdoppeln (Muster wiederholt sich, 8 $\to$ 16 Hex-Ziffern), (2) ein Schritt mehr (Maske \code{0x00000000ffffffff}, Shift 32). Schritte $=\log_2 n$: 5 bei 32, 6 bei 64 Bit.}
\fallstrick{$n$ = \textbf{Bits im Wort}, nicht Array-Länge! $\cdot$ O(1) $\neq$ 1 Befehl (nur: Befehlszahl unabh.\ von $n$) $\cdot$ Schleife über Bits = O($n$), Masken-Kaskade = O($\log n$).}
\bereichende

\bereich[cBits]{Multiplikation \& Potenzen über Bits}
\begin{formelreihe}
\formelblock{Shift-and-Add-Mult.\ (O($n$))}{\begin{aligned}[t]&\code{while(b)\{ if(b\&1) r+=a;}\\ &\code{\ a<{}<=1; b>{}>=1; \}}\end{aligned}}
&
\formelblock{Binary Exponentiation (fxtbook §28.5)}{\begin{aligned}[t]&a^e=\textstyle\prod_{e_i=1} a^{2^i},\ \sim\!\log_2 e\ \text{Mult.}\\ &\code{if(e\&1) t*=s; s*=s; e>{}>=1;}\end{aligned}}
&
\beispielblock{$a^{13}$, $13{=}1101_2$}{\begin{aligned}[t]&t{:}\ a\to a\to a^5\to a^{13}\\ &6\ \text{Mult.\ statt}\ 12\end{aligned}}
\end{formelreihe}
\bereichende

% ============================================================
%  Teil 3 — Code-Optimierung (Floating Point nach Geschwindigkeit)
% ============================================================
\teilkopf{3}{Code-Optimierung}

\bereich[cOpt]{Befehlskosten \& CPU}
\begin{formelreihe}
\formelblock{Kostenklassen (Latency)}{\begin{aligned}[t]&\text{1 Cycle: XOR, ADD, SUB, AND, SHIFT}\\ &\text{wenige: MUL}\\ &\textbf{50\text{--}300: DIV, MOD,}\\ &\textbf{pow, exp, log, sin, cos}\end{aligned}}
&
\formelblock{Latency vs.\ Throughput}{\begin{aligned}[t]&\text{Latency = Zeit \emph{ein} Befehl [Cycles]}\\ &\text{Throughput = fertige Befehle/Cycle}\end{aligned}}
&
\beispielblock{superskalar}{\begin{aligned}[t]&\text{ADD: Lat.\ 1, Thr.\ 4/Cycle}\\ &\Rightarrow\ 4\ \text{Addierer parallel}\end{aligned}}
\end{formelreihe}
\infoblock{Superskalar-Tricks (hidden logic): mehrere Execution Units $\cdot$ Register Renaming (mehr phys.\ als log.\ Reg.) $\cdot$ out-of-order execution $\cdot$ speculative execution \& branch prediction. Real: 5--7 Befehle/Cycle.}
\bereichende

\bereich[cOpt]{Optimier-Repertoire}
\begin{formelreihe}
\formelblock{pow ersetzen}{\begin{aligned}[t]&\code{pow(x,2)}\to\code{x*x}\\ &\code{pow(x,3)}\to\code{x*x*x}\\ &\code{pow(x,8)}\to 3\times\ \text{Quadrieren}\end{aligned}}
&
\formelblock{Division $\to$ Kehrwert-Mult.}{\begin{aligned}[t]&\code{inv = 1.0/d;}\ \text{(1$\times$ teuer)}\\ &\code{a[i]*inv}\ \forall i\ \text{(billig)}\end{aligned}}
&
\beispielblock{warum pow teuer}{\code{pow}=e^{y\ln x}\ \text{(transzendent)}}
\\[0.4em]
\formelblock{Loop Unrolling, 4 Akkus}{\begin{aligned}[t]&\code{s0+=A[j]; s1+=A[j+1];}\\ &\code{s2+=A[j+2]; s3+=A[j+3];}\end{aligned}}
&
\formelblock{Branchless (fxtbook §1.11)}{\begin{aligned}[t]&\code{if((unsigned)x > m)}\ \text{statt}\ x{<}0\,\|\,x{>}m\\ &\code{max0: x \& \textasciitilde(x>{}>(BPL-1))}\end{aligned}}
&
\beispielblock{warum 4 Akkus}{\begin{aligned}[t]&\text{1 Summe: Datenabhängigkeit blockiert}\\ &\text{4 unabh.\ $\Rightarrow$ parallel + L1-Cache}\end{aligned}}
\end{formelreihe}
\fallstrick{\code{==} auf double fast immer falsch $\to$ \code{abs(a-b)<=eps} $\cdot$ Branch-Fehlvorhersage verwirft Pipeline $\cdot$ MUL ist billig, DIV/MOD nicht — Modulo in Schleifen vermeiden.}
\bereichende

% ============================================================
%  Teil 4 — Asymptotik & Relationen
% ============================================================
\teilkopf{4}{Asymptotik \& Relationen}

\bereich[cRel]{Asymptotik}
\begin{formelreihe}
\formelblock{O-Definition}{T(n)\leq c\cdot f(n)\ \ \forall n\geq n_0,\ \ c>0}
&
\formelblock{bei Bit Wizardry}{n=\text{Bits im Wort (nicht Datenmenge!)}}
&
\beispielblock{Klassen erkennen}{\begin{aligned}[t]&\text{Schleife über Bits}\to O(n)\\ &\text{halbierende Masken}\to O(\log n)\\ &\text{feste Befehlszahl}\to O(1)\end{aligned}}
\end{formelreihe}
\infoblock{„O(1) hei\ss t nicht nur 1 Befehl``: Befehlszahl konstant \& unabh.\ von $n$ — 3 oder 12 Befehle sind auch O(1). $\Omega$/$\Theta$ analog mit $\geq$ / beidseitig (s.\ INF4/1).}
\bereichende

\bereich[cRel]{Relationen — „zeige: keine ÄR``}
\begin{formelreihe}
\formelblock{Äquivalenzrelation $\Leftrightarrow$ R+S+T}{\begin{aligned}[t]&\text{R: } a\sim a\ \forall a\\ &\text{S: } a\sim b\Rightarrow b\sim a\\ &\text{T: } a\sim b\ \&\ b\sim c\Rightarrow a\sim c\end{aligned}}
&
\formelblock{Widerlegungs-Strategie}{\begin{aligned}[t]&\text{\textbf{ein} Gegenbeispiel für}\\ &\text{\textbf{eine} Eigenschaft genügt}\end{aligned}}
&
\beispielblock{eps-Vergleich nicht transitiv}{\begin{aligned}[t]&a\sim b :\Leftrightarrow |a-b|\leq\varepsilon,\ \varepsilon{=}1{:}\\ &0\sim 1,\ 1\sim 2,\ \text{aber } |0-2|{=}2{>}1\ \Rightarrow\ 0\nsim 2\end{aligned}}
\end{formelreihe}
\begin{rezept}
\rs{R prüfen: $a\sim a$? (Gegenbsp.\ oft $a=0$)}
\rs{S prüfen: Ordnungsrelationen ($\leq$, $<$, teilt) scheitern hier}
\rs{T prüfen: Toleranz-/Nachbarschaftsrelationen scheitern hier (Fehler addieren sich)}
\rs{verletzte Eigenschaft + konkretes Zahlen-Gegenbeispiel hinschreiben — fertig}
\end{rezept}
\infoblock{Positivbeispiele: $a\equiv b \pmod m$ ($m\,|\,a{-}b$) ist ÄR mit Klassen $\{0,\dots,m{-}1\}$ $\cdot$ „$a{-}b$ gerade`` = mod 2 $\cdot$ ÄR zerlegt Menge in Äquivalenzklassen. $\cdot$ Auf $n$-elementiger Menge: $2^{n^2}$ Relationen.}
\fallstrick{R \& S allein reichen nicht (eps-Bsp.\ erfüllt beide!) $\cdot$ $a\cdot b>0$ scheitert an R wegen $a{=}0$ $\cdot$ $\leq$ scheitert an S, nicht an R/T.}
\bereichende


\end{document}
